\documentclass[aspectratio=169,xcolor={dvipsnames,table}]{beamer}
\usepackage[no-math,deluxe,haranoaji]{luatexja-preset}
\renewcommand{\kanjifamilydefault}{\gtdefault}
\renewcommand{\emph}[1]{{\upshape\bfseries #1}}
\usetheme{metropolis}
\metroset{block=fill}
\setbeamertemplate{navigation symbols}{}
\setbeamertemplate{blocks}[rounded][shadow=false]
\usecolortheme[rgb={0.7,0.2,0.2}]{structure}
%%%%%%%%%%%%%%%%%%%%%%%%%%
%% Change alert block colors
%%% 1- Block title (background and text)
\setbeamercolor{block title alerted}{fg=mDarkTeal, bg=mLightBrown!45!yellow!45}
\setbeamercolor{block title example}{fg=magenta!10!black, bg=mLightGreen!70}
%%% 2- Block body (background)
\setbeamercolor{block body alerted}{bg=mLightBrown!25}
\setbeamercolor{block body example}{bg=mLightGreen!15}
%%%%%%%%%%%%%%%%%%%%%%%%%%%
\usepackage[absolute,overlay]{textpos}  %% 任意の位置に図を配置
%\usepackage[grid=true,gridcolor=Maroon,subgridcolor=gray,gridunit=pt,texcoord]{eso-pic} %%場所決めのための格子 \myPreamble4Slide.texより先に読み込むこと
%%%%%%%%%%%%%%%%%%%%%%%%%%%
%%%%%%%%%%%%%%%%%%%%%%%%%%%
%% さまざまなアイコン
%%%%%%%%%%%%%%%%%%%%%%%%%%%
%\usepackage{fontawesome}
\usepackage{fontawesome5}
\usepackage{figchild}
\usepackage{twemojis}
\usepackage{utfsym}
\usepackage{bclogo}
\usepackage{marvosym}
\usepackage{fontmfizz}
\usepackage{pifont}
\usepackage{phaistos}
\usepackage{worldflags}
\usepackage{jigsaw}
\usepackage{tikzlings}
\usepackage{tikzducks}
\usepackage{scsnowman}
\usepackage{epsdice}
\usepackage{halloweenmath}
\usepackage{svrsymbols}
\usepackage{countriesofeurope}
\usepackage{tipa}
\usepackage{manfnt}
%%%%%%%%%%%%%%%%%%%%%%%%%%%
\usepackage{tikz}
\usetikzlibrary{calc,patterns,decorations.pathmorphing,backgrounds}
\usepackage{tcolorbox}
\usepackage{tikzpeople}
\usepackage{circledsteps}
\usepackage{xcolor}
\usepackage{amsmath}
\usepackage{booktabs}
\usepackage{chronology}
\usepackage{signchart}
%%%%%%%%%%%%%%%%%%%%%%%%%%%
%% 場合分け
%%%%%%%%%%%%%%%%%%%%%%%%%%%
\usepackage{cases}
%%%%%%%%%%%%%%%%%%%%%%%%%%
\usepackage{pdfpages}
%%%%%%%%%%%%%%%%%%%%%%%%%%%
%% 音声リンク表示
\newcommand{\myaudio}[1]{\href{#1}{\faVolumeUp}}
%%%%%%%%%%%%%%%%%%%%%%%%%%
%% \myAnch{<名前>}{<色>}{<テキスト>}
%% 指定のテキストを指定の色の四角枠で囲み, 指定の名前をもつTikZの
%% ノードとして出力する. 図には remember picture 属性を付けている
%% ので外部から参照可能である.
\newcommand*{\myAnch}[3]{%
  \tikz[remember picture,baseline=(#1.base)]
    \node[draw,rectangle,line width=1pt,#2] (#1) {\normalcolor #3};
}
%%%%%%%%%%%%%%%%%%%%%%%%%%
%% \myEmph コマンドの定義
%%%%%%%%%%%%%%%%%%%%%%%%%%
%\newcommand{\myEmph}[3]{%
%    \textbf<#1>{\color<#1>{#2}{#3}}%
%}
\usepackage{xparse} % xparseパッケージの読み込み
\NewDocumentCommand{\myEmph}{O{} m m}{%
    \def\argOne{#1}%
    \ifx\argOne\empty
        \textbf{\color{#2}{#3}}% オプション引数が省略された場合
    \else
        \textbf<#1>{\color<#1>{#2}{#3}}% オプション引数が指定された場合
    \fi
}
%%%%%%%%%%%%%%%%%%%%%%%%%%%
%%%%%%%%%%%%%%%%%%%%%%%%%%%
%% 文末の上昇イントネーション記号 \myRisingPitch
%% 通常のイントネーション \myDownwardPitch
%% https://note.com/dan_oyama/n/n8be58e8797b2
%%%%%%%%%%%%%%%%%%%%%%%%%%%
\newcommand{\myRisingPitch}{
\begin{tikzpicture}[scale=0.3,baseline=0.3]
\draw[->,>=stealth] (0,0) to[bend right=45] (1,1);
\end{tikzpicture}
}
\newcommand{\myDownwardPitch}{
\begin{tikzpicture}[scale=0.3,baseline=0.3]
\draw[->,>=stealth] (0,1) to[bend left=45] (1,0);
\end{tikzpicture}
}
%%%%%%%%%%%%%%%%%%%%%%%%%%%%
%\AtBeginSection[%
%]{%
%  \begin{frame}[plain]\frametitle{授業の流れ}
%     \tableofcontents[currentsection]
%   \end{frame}%
%}

\usetikzlibrary{tikzmark}
\usepackage{pxrubrica}
\usepackage{pangram}
\PangramSetup{
sizes={7pt,8pt,9pt,10pt,11pt,12pt,13pt,14pt,15pt,16pt},
tagfont=\color{gray}\ttfamily\footnotesize,}
\NewPangramClass{IRO}{\input{iroha.txt}}
%%%%%%%%%%%%%%%%%%%%%%%%%%
%%%%%%%%%%%%%%%%%%%%%%%%%%%
\title{English is fun.}
\subtitle{The quick brown fox jumps over the lazy dog.}
\author{{\tiny 0054}\,{\scriptsize\myaudio{./audio/020_part_of_speech_00.mp3}}
}
\institute[]{}
\date[]

%%%%%%%%%%%%%%%%%%%%%%%%%%%%
%% TEXT
%%%%%%%%%%%%%%%%%%%%%%%%%%%%
\begin{document}

\begin{frame}[plain]
  \titlepage
\end{frame}

\section*{授業の流れ}
\begin{frame}[plain]
  \frametitle{授業の流れ}
  \tableofcontents
\end{frame}

%%%%%%%%%%%%%%%%%%%%%%%%%%%%%%%%%%%%%%%%%%%%
\section{品詞}
\subsection{品詞とは}
%%%%%%%%%%%%%%%%%%%%%%%%%%%%%%%%%%%%%%%%%%%%%
\begin{frame}[plain,label=basic]{品詞}
\large 
\begin{description}
 \item<2->[名詞] 人、動物、植物やモノの名前を表す語\hfill\visible<3->{\small child, cat, flower, desk \ldots}
  \item<4->[代名詞] 名詞の代わりに使う語\hfill\visible<5->{\small I, we, you he, she, it, they, this, that \ldots}
 \item<6->[動詞] 動作や状態を表す語\hfill\visible<7->{\small run, see, speak, eat, study, have, be \ldots{}}
 \item<8->[形容詞] 名詞をくわしく説明する語\hfill\visible<9->{\small big, new, happy, good \ldots}
 \item<10->[副詞] 動詞や形容詞などをくわしく説明する語\hfill\visible<11->{\small quickly, here, then\ldots}
 \item<12->[前置詞] 名詞や代名詞と組み合わせて$+\alpha$の意味をくわえる\\
\hfill\visible<13->{\small in, on, under, with, before, after, from, to \ldots}
\end{description}

\vfill

\hfill\visible<14->{まずはこれをしっかり理解しましょう}
\end{frame}
%%%%%%%%%%%%%%%%%%%%%%%%%%%%%%%%%%%%%%%%%%%%%%
\begin{frame}[plain]{Exercises}
それぞれの英文の単語の品詞がなにかかんがえましょう%
\hfill{\tiny 0232}\,{\scriptsize \myaudio{./audio/020_part_of_speech_01.mp3}}

\begin{columns}[t]
\begin{column}{.7\textwidth}
  \begin{enumerate}
  \item Small cats jump quickly.\\
	\visible<2->{{\small \Circled{ 形 }\hspace{12pt}\Circled{ 名 } \Circled{ 動 }\hspace{12pt}\Circled{ 副 }}}
  \item Tall trees grow slowly.\\
	\visible<3->{{\small \Circled{ 形 } \Circled{ 名 }\hspace{8pt}\Circled{ 動 }\hspace{10pt}\Circled{ 副 }}}
  \item Big dogs bark loudly.\\
	\visible<4->{{\small \Circled{ 形 } \Circled{ 名 } \Circled{ 動 }\hspace{10pt}\Circled{ 副 }}}
  \item George had very old coins.\\
	\visible<5->{{\small \Circled{ 名 }\hspace{18pt}\Circled{ 動 } \Circled{ 副 } \Circled{ 形 }} \Circled{ 名 }}
  \item Jennifer was very happy then.\\
	\visible<6->{{\small \Circled{ 名 }\hspace{22pt}\Circled{ 動 } \Circled{ 副 }\hspace{8pt}\Circled{ 形 }\hspace{10pt}\Circled{ 副 }}}
 \end{enumerate}
\end{column}
\begin{column}{.25\textwidth}
 \begin{tcolorbox}
  \begin{tabular}{cl}
    \Circled{ 名 }&名詞\\
    \Circled{ 動 }&動詞\\
    \Circled{ 形 }&形容詞\\
    \Circled{ 副 }&副詞
  \end{tabular}
 \end{tcolorbox}
\end{column}
\end{columns}
\end{frame}
%%%%%%%%%%%%%%%%%%%%%%%%%%%%%%%%
\begin{frame}[plain]{代名詞}
\large

代名詞: \kenten{名}詞の\kenten{代}わりに使う語

 \begin{enumerate} 
  \item \begin{enumerate}
	 \item<1-> Dave is my classmate.\hspace{5pt} Dave likes basketball.\\
	      \mbox{}
	 \item<2-> \myAnch{X}{white}{Dave}is my classmate.\hspace{10pt}\myAnch{x}{Maroon}{He} likes basketball.\\
	      \mbox{}
	\end{enumerate}  \item \begin{enumerate}
	 \item<4-> I ate an apple.\hspace{5pt} The apple was delicious.%
	      \hfill{\scriptsize ate \textipa{/\'eit/} 食べた(eatの過去形です)}\\
\hfill{\scriptsize delicious \textipa{/dIl\'IS@s/} おいしい}\\
%	      \mbox{}
	 \item<5-> I ate \myAnch{Y}{white}{an apple.}\hspace{10pt} \myAnch{y}{Maroon}{It} was delicious.
	\end{enumerate}
 \end{enumerate}

\begin{tikzpicture}[remember picture, overlay]
% Calculate intermediate points
  \coordinate (X1) at ($(X) + (0,15pt)$); % 10pt above X
  \coordinate (x1) at ($(x) + (0,15pt)$); % 10pt above x
  % Draw the arrow with right angles
  \visible<3->{\draw[<-,Maroon] (X) -- (X1) -- (x1) -- (x);}
%%%%%%%%%%%
% Calculate intermediate points
  \coordinate (Y1) at ($(Y) + (0,15pt)$); % 10pt above Y
  \coordinate (y1) at ($(y) + (0,15pt)$); % 10pt above y
  % Draw the arrow with right angles
  \visible<6->{\draw[<-,Maroon] (Y) -- (Y1) -- (y1) -- (y);}
\end{tikzpicture}
\hfill{\tiny 0309}\,{\scriptsize \myaudio{./audio/020_part_of_speech_03.mp3}}
\end{frame}
%%%%%%%%%%%%%%%%%%%%%%%%
%%%%%%%%%%%%%%%%%%%%%%%%%%%%%%%
\begin{frame}[plain]{前置詞}
\hfill{\tiny 0234}\,{\scriptsize \myaudio{../1st_grader/audio/004_verb_object_4.mp3}}
 \begin{enumerate}
  \item<1-> He plays tennis.
  \item<2-> He plays tennis \Circled[fill color = gray, outer color=gray, inner color=white]{\,\,with his friends\,\,}.%
\hfill{\scriptsize with: ～といっしょに}
  \item<3-> He plays tennis \Circled[fill color = gray, outer color=gray, inner color=white]{\,\,in the park\,\,}.%
\hfill{\scriptsize in: ～の中で}
  \item<4-> He plays tennis \Circled[fill color = gray, outer color=gray, inner color=white]{\,\,after school\,\,}.%
\hfill{\scriptsize after: ～の後に}
  \item<5-> He plays tennis \Circled[fill color = gray, outer color=gray, inner color=white]{\,\,with his friends\,\,} \Circled[fill color = gray, outer color=gray, inner color=white]{\,\,in the park\,\,} \Circled[fill color = gray, outer color=gray, inner color=white]{\,\,after school\,\,}.

 \end{enumerate}

\begin{block}<6->{Topics for Today}\small
with, in, after などを\kenten{前置詞}といいます%
\hfill{}「（代）名詞の\kenten{前}に\kenten{置}く\kenten{詞}」
\begin{itemize}\setbeamertemplate{items}[square]\small
 \item<7-> 裏返していうと「前置詞の\kenten{後}にはかならず（代）\kenten{名詞}がくる」ということです
 \item<8-> \Circled[fill color = gray, outer color=gray, inner color=white]{\,\,$\text{前置詞}+\text{（代）名詞}$\,\,}\,は、文に$+\alpha$の意味をつけくわえます\\
 \hfill\visible<9->{{\scriptsize 前置詞はいろいろあります。ちょっとずつ覚えましょう}}
\end{itemize}
     \end{block}
\end{frame}
%%%%%%%%%%%%%%%%%%%
\againframe<14>[plain]{basic}
%%%%%%%%%%%%%%%%%%%%%%%%
%\begin{frame}[plain]{その他}
% \begin{description}
%  \item[冠詞]  名詞の前に置く\hfill{}a / an, theだけ
%  \item[助動詞] 動詞の意味を補う語\hfill{}can, will, must \ldots 
% \end{description}
%\end{frame}
%%%%%%%%%%%%%%%%%%%%%%
\section{冠詞}
\begin{frame}[plain]{冠詞}
\large

冠詞: a / an ,\hspace{10pt}the

 \begin{enumerate}
  \item $\text{a / an} + \text{単数形の名詞}$\,(1つの～、ある～)
	\begin{enumerate}
	 \item a book\hfill\textipa{/\textschwa /}
	 \item an apple\hfill\textipa{/\textschwa n/}
	\end{enumerate}
  \item $\text{the} + \text{名詞}$\,(その～)\hspace{10pt}{\small (単数形・複数形どちらにもつく)}
	\begin{enumerate}
	 \item the book / the books\hfill\textipa{/\dh\textschwa/}
	 \item the apple / the apples\hfill\textipa{/\dh i/}
	\end{enumerate}
 \end{enumerate}
\hfill{\tiny 0218}\,{\scriptsize \myaudio{./audio/020_part_of_speech_02.mp3}}
\end{frame}
%%%%%%%%%%%%%%%%%%%%%%%%%%%%
\begin{frame}[plain]{冠詞}
 
\begin{enumerate}
 \item She has \textbf{a} cat.\hspace{5pt} \textbf{The} cat is black.%
 \item I ate \textbf{an} apple.\hspace{5pt} \textbf{The} apple was delicious.%
	      \hfill{\scriptsize ate \textipa{/\'eIt/} 食べた(eatの過去形です)}\\
\hfill{\scriptsize delicious \textipa{/dIl\'IS@s/} おいしい}
\end{enumerate}

\vfill

\begin{block}{Topics for Today}
 $\text{冠詞} = \left\{\begin{tabular}{ll}a(an)&不定冠詞\\the&定冠詞\end{tabular}\right.$
\end{block}

{\tiny 0144}\,\hfill{\scriptsize \myaudio{./audio/020_part_of_speech_99.mp3}}

\end{frame}
%%%%%%%%%%%%%%%%%%%%%%%%%%%%%
%\begin{frame}[plain]{助動詞}
% \Large
%
%助動詞: 動詞の意味を補う
%
%\begin{enumerate}
% \item \begin{enumerate}
%	\item<1-> I play the guitar.
%	\item<2-> I \myEmph[2-]{Maroon}{will} play the guitar tomorrow.（未来）
%       \end{enumerate}
% \item \begin{enumerate}
%	\item<3-> You play the guitar.
%	\item<4-> You \myEmph[3-]{Maroon}{must} play the guitar.（義務）
%       \end{enumerate}
% \item \begin{enumerate}
%	\item<5-> He plays the guitar.
%	\item<6-> He \myEmph[4-]{Maroon}{can} play the guitar.（可能）
%       \end{enumerate}
%\end{enumerate}
%\mbox{}\hfill{\scriptsize \myaudio{./audio/020_part_of_speech_04.mp3}}
%\end{frame}
%%%%%%%%%%%%%%%%%%%%%%%%
\section{品詞のまとめ}
%%%%%%%%%%%%%%%%%%%%%%%
\againframe{basic}
%%%%%%%%%%%%%%%%%%%%%%%
\section{混乱しないで}
%%%%%%%%%%%%%%%%%%%%%%
\begin{frame}[plain]{きちんと区別しよう}
 \begin{enumerate}
  \item<1-> birdの品詞は？\hfill\visible<2->{名詞}
  \item<3-> Birds sing.という英文のBirdsの品詞は？\hfill\visible<4->{名詞}
  \item<5-> Birds sing. という英文の主語は？\hfill\visible<6->{Birds}
 \end{enumerate}

\visible<7->{てことはBirds sing.のBirdsは名詞でもあるし主語でもあるってことなの？}

\visible<8->{Birds sing.のBirdsは「名詞です」といったり「主語です」といったり、なんか混乱しちゃうんですけど}
\end{frame}
%%%%%%%%%%%%%%%%%%%%%%
\begin{frame}[plain]{きちんと区別しよう}

説明の角度がちがいます\pause

\bigskip

「Birds sing.のBirdsは主語です」というときは、Birdsという単語が文の中でどういう役割をしているかについて言及しています。\pause
おなじbirdsでもI love birds.という文のbirdsは目的語ですよね\pause

\bigskip

いっぽう「Birds sing.のBirdsは名詞です」というときは、文中での役割とは無関係にbirdsという単語固有の話をしています。いいかえるとbirdsという単語は文のどこに出てきても、いつでも名詞なのです。
I love birds.のbirdsが名詞なのは明らかですよね
\end{frame}
%%%%%%%%%%%%%%%%%%%%%%%%
\begin{frame}[plain]{きちんと区別しよう}

\visible<1->{{\Large We drink \textcolor{Maroon}{\bfseries water}\tikzmark{water1}. \textcolor{Maroon}{\bfseries Water}\tikzmark{water2} is very important.}}

\vfill

\hfill\hfill{}\tikzmark{noun}\visible<2->{名詞}\hfill~

\vfill

\hspace*{100pt}\visible<3->{目的語}\tikzmark{object}\hspace{60pt}\visible<4->{主語}\tikzmark{subject}

\begin{tikzpicture}[remember picture,overlay]
 \visible<2->{\draw[<-,Maroon,opacity=.7,line width=2pt] ([xshift=-15pt,yshift=-2pt]pic cs:water1) to[out=-70, in=180] ([xshift=-2pt, yshift=4pt] pic cs:noun);}
 \visible<2->{\draw[<-,Maroon,opacity=.7,line width=2pt] ([xshift=-14pt,yshift=-2pt]pic cs:water2) to[out=-70, in=180] ([xshift=-2pt, yshift=4pt] pic cs:noun);}
%%%%%%%%%%%%%%%
 \visible<3->{\draw[->,opacity=.4,line width=2pt] ([xshift=-15pt,yshift=-2pt]pic cs:water1) to[out=-90, in=90] ([xshift=-14pt, yshift=10pt] pic cs:object);}
 \visible<4->{\draw[->,opacity=.4,line width=2pt] ([xshift=-15pt,yshift=-2pt]pic cs:water2) to[out=-90, in=90] ([xshift=-14pt, yshift=10pt] pic cs:subject);}
\end{tikzpicture}

\begin{textblock*}{0.4\linewidth}(275pt,165pt)
    % TiKZを使った図形の描画
    \visible<5->{\begin{tikzpicture}
      \node[name=k, shape=businessman, minimum size=55pt, mirrored] {};
  \node[ellipse callout, xshift=-1.8cm, draw,yshift= .2cm,
    callout absolute pointer={(k.mouth)},
    font=\tiny\bfseries, align=center, inner sep=1pt]
    {名詞は\\SにもOにも\\なります};
    \end{tikzpicture}}
\end{textblock*}
\end{frame}
%%%%%%%%%%%%%%%%%%%%%%%%%
\section{おまけ}
%%%%%%%%%%%%%%%%%%%%%%%%%
\begin{frame}[plain]{???}
\Large
 The qui\myEmph[5-]{Maroon}{c}k \myEmph[4-]{Maroon}{b}rown fox jumps over the l\myEmph[3-]{Maroon}{a}zy \myEmph[6-]{Maroon}{d}og. 

\visible<2->{
\small \Circled{ 冠 }\hspace{12pt}
\small \Circled{ 形 }\hspace{17pt}
\small \Circled{ 形 }\hspace{16pt}
\small \Circled{ 名 }\hspace{12pt}
\small \Circled{ 動 }\hspace{14pt}
\small \Circled{ 前 }\hspace{8pt}
\small \Circled{ 冠 }\hspace{5pt}
\small \Circled{ 形 }\hspace{7pt}
\small \Circled{ 名 }

}
\vfill

\visible<7->{上の英文から\\小文字のアルファベットをaから順番にさがしてみよう}

\visible<8->{すべての文字を使用した文$\rightarrow$パングラム}
\end{frame}
%%%%%%%%%%%%%%%%%%%%%%%%%%
\begin{frame}[plain]{フォントのサンプル}
 \pangram
\end{frame}
%%%%%%%%%%%%%%%%%%%%%%%%%%%
\begin{frame}[plain]{日本語のパングラム}
\pause
 \pangram[textclass=IRO,font=\fontspec{Latin Modern Mono}]  % いろは歌設定
\end{frame}
%%%%%%%%%%%%%%%%%%%%%%%%%%
\end{document}

